%% kapa: exam=1 round=34 year=2023 period=1 form=A subject=civil_law subject_ko=민법 seq=18 number=18
%% source: Q-Net 감정평가사 기출문제 (cst00309) artlSeq=5211789 file=raw/1008/5211789/2023년도 제34회 감정평가사 1차시험 1교시.pdf
%% conversion: kordoc 4.24.51 -> tools/gen1.py (auto); review=proofread-sonnet; flags=none
\begin{question}{18}
법률행위 부관인 조건에 관한 설명으로 옳지 않은 것은? (다툼이 있으면 판례에 따름)
\begin{choices}
  \item 물권행위에는 조건을 붙일 수 없다.
  \item 조건이 되기 위해서는 법률이 요구하는 것이 아니라 당사자가 임의로 부가한 것이어야 한다.
  \item 조건의 성취를 의제하는 효과를 발생시키는 조건성취 방해행위에는 과실에 의한 행위도 포함된다.
  \item 부첩(夫妾)관계의 종료를 해제조건으로 하는 부동산 증여계약은 해제조건뿐만 아니라 증여계약도 무효이다.
  \item 당사자의 특별한 의사표시가 없는 한 정지조건이든 해제조건이든 그 성취의 효력은 소급하지 않는다.
\end{choices}
\end{question}
\answer{1}
